\documentclass[nooutcomes,noauthor, handout]{ximera}
%\documentclass{ximera}
\input{../../preamble}

\title{Area and Volume Conversions}
\begin{document}
\begin{abstract}\end{abstract}
\maketitle



\begin{problem}
	Jasmine has a bunch of Post-its which are is $3$-inch squares. She has covered a poster with an $8$-by-$14$ array of these Post-its. 
	\begin{enumerate}
		\item What is the area of the poster in Post-its? Explain any calculations you do.
		\item What is the area of the poster in square inches? Explain any calculations you do.
	\end{enumerate}
\end{problem}



\begin{problem}
Aurora says, ``There must be $3$ square feet in a square yard, because there are $3$ feet in one yard.''

\begin{enumerate}
	\item Draw pictures to help Aurora investigate this statement. How many square feet are in a square yard?
	\item How is this problem related to Jasmine's problem about Post-its? Explain your thinking.
\end{enumerate}
\end{problem}



\begin{problem}
Tiana has packed a shipping box full of smaller boxes. The smaller boxes are $10$cm long, $7$cm wide, and $15$cm tall. The shipping box fits $120$ of these smaller boxes.
\begin{enumerate}
	\item What is the volume of the shipping box in cubic centimeters? Explain any calculations you do.
	\item Did you need to know the shape of the shipping box to solve this problem? Why or why not? Explain your thinking.
\end{enumerate}
\end{problem}



\begin{problem}
Merida says, ``I think $1$ cubic meter will be $100$ cubic centimeters.''

\begin{enumerate}
	\item Explain what Merida is thinking, and sketch pictures to illustrate her ideas. How many cubic centimeters are in a cubic meter?
	\item How is this problem related to Tiana's problem about boxes? Explain your thinking.
\end{enumerate}
\end{problem}




\begin{problem}
	Using the fact that there are $2.54$ centimeters in every inch, convert $25$ square inches into square centimeters. 
	
	You should solve this problem in two different ways, drawing pictures as needed to illustrate your thinking.
\end{problem}




\begin{problem}
	Using the fact that there are $2.54$ centimeters in every inch, convert $25$ square centimeters into square inches. 
	
	You should solve this problem in two different ways, drawing pictures as needed to illustrate your thinking.
\end{problem}











\newpage

\begin{instructorNotes}

{\bf Main goal:} We use the Math 1125 meanings of multiplication and division to convert between units using area and volume.

{\bf Overall picture:} 

This activity is designed to help students build skills for converting area and volume between unusual units, within the English and metric system, and then between English and metric units.  As with other conversion activities, we are not looking for dimensional analysis or making a ``fence post'' diagram. Dimensional analysis tends to be another ``rule'' that students have memorized, and obscures the meaning of multiplication and division in conversion problems.


\begin{itemize}
\item Two nice solution methods to see here are (a) to draw a rectangle or prism with the given volume, convert each of the lengths of the sides, and then use the converted lengths plus an area or volume formula to find the total or (b) to convert a single unit of area or volume, and then multiply by the total number of units you need.

\item Students also don't naturally try to give their answers in terms of the meaning of multiplication and division, and will need to be coaxed in this direction.

	\item Encourage the students to try to think about this problem in multiple ways.   Highlight all different solutions you see, as well as any misconceptions that students are willing to share. In particular, watch out for the idea that if the length is multiplied by $k$, the area should also be multiplied by $k$. Relate this back to our work with scaling area and volume. 
	\item Encourage students to draw a picture if they're stuck.  Using two different methods could simply involve drawing two different pictures to solve the problem!
	\item If the students draw a picture, have them explain how they know they can draw that particular picture.  Essentially, this is a consequence of moving and additivity.  We can rearrange the area in any fashion we like because of the principles of moving and additivity - no area is created or lost when we rearrange. 
	\item Emphasize the meaning of multiplication and division in these problems.  Students should identify the groups and objects to describe why they decide to multiply, or why they decide to divide.  With division, students should be able to identify which type of division they are using - though in the case of conversion it is usually ``how many groups?''.
\end{itemize}


{\bf Suggested timing: } You can give students 20 mintues to work on these problems, then spend the rest of the time in presentations and discussion. Or, you could split the work into three phases: problems 1 and 2 about length (5 minutes work, 10 minutes discuss), problems 3 and 4 about volume (5 minutes work, 10 minutes discuss), and finally problems 5 and 6 converting between English and metric (5 minutes work, 10 minutes discuss).

\end{instructorNotes}



\end{document}